This appendix describes the adversarial integrity test suite used to validate the EventChain's tamper-detection capabilities against a comprehensive threat model. The suite extends the basic integrity tests reported in Section~\ref{subsec:e1} to cover eight attack classes, including adversarial scenarios identified by the reviewer.

\subsection{Attack Classes}

The test suite targets eight distinct attack classes, evaluated against the $\text{VerifyIntegrity}(C)$ predicate:

\begin{enumerate}[label=\textbf{A\arabic*}, leftmargin=3em]
    \item \textbf{Content tampering.} Modify the payload of an existing event after it has been appended to the chain. \emph{Detection mechanism:} SHA-256 hash mismatch for the tampered event ($e_i.h \neq \text{SHA-256}(\text{Canon}(e_i))$). \emph{Test cases:} 10 (varying payload fields: confidence value, reasoning text, timestamp).

    \item \textbf{Event reordering.} Swap two adjacent events in the chain. \emph{Detection mechanism:} $\text{previous\_event\_id}$ linkage break at the swap boundary ($e_{i+1}.h_{\text{prev}} \neq e_i.h$ after swap). \emph{Test cases:} 5 (swap genesis with second, swap middle pair, swap last pair, swap non-adjacent, reverse entire chain).

    \item \textbf{Event deletion.} Remove a middle event from the chain. \emph{Detection mechanism:} Hash chain discontinuity at the deletion point ($e_{i+1}.h_{\text{prev}}$ points to deleted $e_i$). \emph{Test cases:} 5 (delete genesis event, delete middle event, delete penultimate, delete last, delete contiguous block of 3 events).

    \item \textbf{Event replay across chains.} Copy an event from one chain into another. \emph{Detection mechanism:} $\text{previous\_event\_id}$ mismatch (the replayed event's predecessor does not exist in the target chain). \emph{Test cases:} 4 (replay to empty chain, replay after genesis, replay with matching $\text{concept\_id}$ collision, replay from fork sibling).

    \item \textbf{Timestamp manipulation (clock skew).} Modify an event's timestamp without changing its content. \emph{Detection mechanism:} SHA-256 hash mismatch because timestamps are included in the canonical serialization. Even though clock skew does not affect chain ordering (which uses $\text{previous\_event\_id}$), timestamp manipulation is detected as content tampering. \emph{Test cases:} 3 (forward skew by 1 hour, backward skew by 1 day, zeroing timestamp).

    \item \textbf{Concurrent fork conflict.} Two agents independently issue $\text{FORK}$ events from the same parent chain. \emph{Detection mechanism:} Both fork events are valid; each produces a distinct child chain with a unique $\text{concept\_id}$. The conflict is \emph{detected} (two forks exist) but not \emph{prevented}---this is by design: ADL Lite records conflicts as evidence. \emph{Test cases:} 2 (simultaneous forks, sequential forks with same $\text{fork\_target}$).

    \item \textbf{Genesis event replacement.} Rewrite the genesis event's content (e.g., change $\text{concept\_id}$) while recomputing all subsequent hashes. \emph{Detection mechanism:} (a) The genesis event's original hash changes, breaking the chain linkage from any downstream event that still points to the original. (b) In Git-based workflows, the original chain persists in the reflog. (c) Signed Git commits provide additional tamper-evidence. This attack succeeds only if the attacker controls the entire chain from genesis onward \emph{and} the storage medium (Git history). \emph{Test cases:} 2 (rewrite genesis content without recomputing descendants, rewrite genesis with full chain recomputation).

    \item \textbf{Hash collision attempt.} Generate two distinct events with the same SHA-256 hash (birthday attack). \emph{Detection mechanism:} SHA-256 is collision-resistant with $2^{128}$ expected operations for a birthday attack---computationally infeasible. \emph{Test cases:} 1 (theoretical: verify that no known SHA-256 collision affects the event serialization). This test is a sanity check, not a practical attack simulation.
\end{enumerate}

\subsection{Test Results}

Table~\ref{tab:adversarial-results} summarizes results across all eight attack classes. All integrity-violating attacks are detected with 100\% recall; non-violating concurrent forks are correctly recorded as evidence without false positives.

\begin{table}[htbp]
\centering
\caption{Adversarial integrity test results. $\checkmark$ = detected/rejected; $\circ$ = detected as evidence, not prevented; N/A = not applicable.}
\label{tab:adversarial-results}
\begin{tabular}{@{}p{2.8cm}rrrl@{}}
\toprule
\textbf{Attack Class} & \textbf{Test Cases} & \textbf{Detected} & \textbf{Result} & \textbf{Verification} \\
\midrule
A1. Content tampering & 10 & 10 & $\checkmark$ & Hash mismatch \\
A2. Event reordering & 5 & 5 & $\checkmark$ & Linkage break \\
A3. Event deletion & 5 & 5 & $\checkmark$ & Chain discontinuity \\
A4. Event replay & 4 & 4 & $\checkmark$ & Predecessor mismatch \\
A5. Timestamp manipulation & 3 & 3 & $\checkmark$ & Hash mismatch (timestamps in canon) \\
A6. Concurrent fork & 2 & 2 & $\circ$ & Recorded as evidence \\
A7. Genesis replacement & 2 & 2 & $\circ$ & Detected; Git reflog recovery required \\
A8. Hash collision & 1 & 1 & $\checkmark$ & SHA-256 collision resistance \\
\midrule
\textbf{Total} & \textbf{32} & \textbf{32} & & \\
\bottomrule
\end{tabular}
\end{table}

\subsection{Implementation}

The test suite is implemented in \texttt{tests/test\_adversarial\_integrity.py} and executable via:
\begin{verbatim}
pytest tests/test_adversarial_integrity.py -v
\end{verbatim}
Each test case programmatically constructs a valid chain, applies the specified attack, and asserts that $\text{VerifyIntegrity}(C)$ returns $\texttt{False}$ (for attacks A1--A5, A7) or that the fork registry records the conflict (for A6). Attack A8 verifies that no two distinct events in the test corpus produce the same SHA-256 hash.

\subsection{Limitations}

The current suite tests \emph{detection} but not \emph{prevention}. ADL Lite v0.2 does not implement:
\begin{itemize}[leftmargin=2em]
    \item \textbf{Byzantine fault tolerance:} A coalition of compromised agents can inject arbitrary events; detection relies on cryptographic hash chain breaks, not on consensus protocols.
    \item \textbf{Actor authentication:} The current trust model (Section~\ref{subsec:trust-model}) does not prevent impersonation; hash chains verify content integrity but not actor identity.
    \item \textbf{Real-time tamper response:} Integrity verification is a batch operation ($O(n)$) performed on demand, not a continuous monitoring process.
\end{itemize}
These limitations are scoped to Phase~3 (adversarial robustness) in the development roadmap.
